\section{Mathematical Formulations and Tool Interface}
\label{app:mathematical-details}

\subsection{Parameterized Design Space}
\label{app:design-space}
Table~\ref{tab:design-space} specifies the implemented base-only serial family through its generator domains. The family definition composes the base segment, guide, connector, and payload prototypes with the serial assembly rule. A study selects its permitted subset and freezes the associated task and control configuration.
\begin{table*}[!t]
\centering
\caption{Parameterized design space of the tendon-driven soft-robot family: implemented base-only serial configuration.}
\label{tab:design-space}
\papertablestyle
\begin{tabularx}{\textwidth}{@{}>{\itshape}p{0.14\textwidth}!{\color{TraceGray}\vrule width 0.25pt}p{0.24\textwidth}p{0.10\textwidth}Y@{}}
\toprule
\normalfont\textbf{Category} & \textbf{Parameter} & \textbf{Type} & \textbf{Domain or assembly rule} \\
\midrule
Topology & Segment count; tendon count & Discrete & 1--9 segments; 1--7 tendons; consecutive base-origin spans \\
Geometry & Segment length; section family and dimensions & Mixed & 60--240~mm generic length; circle, tube, ellipse, rectangle; indexed bounds in Appendix~\ref{app:design-space} \\
Material & Young's modulus & Continuous & 2--15~MPa per segment \\
Tendon routing & Radius; guide and anchor station; angular placement & Mixed & Radius 14--23~mm; guide $s\in[0.15,0.75]$; anchor $s\in[0.8,1]$; phase $[-2\pi,2\pi]$; offsets $[-\pi,\pi]$~rad \\
Assembly & Tendon diameter; force rating; component placement & Mixed & Diameter 0.5--1.2~mm; force limit 1--8~N; guides and connectors at serial boundaries, payload at the tip \\
\midrule
Control $(U,\theta)$ & Input sequence; initialization; controller settings & Continuous & Task-specific admissible input; pretension 0--1~N; terminal/holding speed weights 0--1 in the base configuration \\
\bottomrule
\end{tabularx}
\end{table*}


Indexed segment domains override the generic domain: segment~0 has length 80--240~mm, segment~1 has length 60--200~mm, and subsequent segments use 60--240~mm. Circular radius and elliptical semiaxes range from 3 to 12~mm; tube outer and inner radii range from 5 to 12~mm and 1 to 4~mm, respectively; rectangular width and height range from 6 to 18~mm. The compiler enforces positive section dimensions, tube wall thickness, ordered route stations, nonzero spans, owned guide holes, routing clearances, and consistent force and pretension limits. The base-only rule starts each tendon span at the base and routes it through consecutive segments to its anchor.

Density, bending viscosity, connector and payload inertia, and prototype station angles are inherited from the selected assembly in this configuration. Active pretension and controller weights belong to the control configuration; passive fixed-payout conditions are specified by the passive assembly and load model. Numerical cell count controls discretization. Parameter identifiers, units, and effective model bindings are resolved for each compiled candidate.


\subsection{Passive and Multi-condition Formulations}
For a passive fixed-payout robot, the response is determined by its structure, unilateral tendon constraints, and applied load. A quasi-static equilibrium $q^\star_\tau(d,f)$ satisfies
\begin{equation}
\begin{aligned}
\nabla_q\Pi(q^\star_\tau;d,f)
+\sum_j\mu^\star_j\nabla_q c_j(q^\star_\tau,d)&=0,\\
c_j(q^\star_\tau,d)\leq0,\quad
\mu^\star_j&\geq0,\quad
\mu^\star_jc_j(q^\star_\tau,d)=0,
\end{aligned}
\label{eq:passive-problem}
\end{equation}
where $\Pi$ is the potential energy and $c_j$ denotes a unilateral tendon or assembly constraint. The passive response is evaluated against
\begin{equation}
g^{\mathrm{pass}}_{\tau,i}
\bigl(d,q^\star_\tau(d,f)\bigr)\leq0.
\label{eq:passive-constraints}
\end{equation}

For a set of required conditions $\mathcal Q_\tau=\mathcal Q_{\mathrm{act}}\cup\mathcal Q_{\mathrm{pass}}$, one morphology is shared while each active condition may use its own control sequence. The condition-wise violation is
\begin{equation}
v_c(d)=
\begin{cases}
\displaystyle\inf_{U_c\in\mathcal U_c(d)}
\max_i g_{c,i}(d,U_c), & c\in\mathcal Q_{\mathrm{act}},\\[0.8em]
\displaystyle\max_i g^{\mathrm{pass}}_{c,i}
\bigl(d,q^\star_c(d,f_c)\bigr), & c\in\mathcal Q_{\mathrm{pass}},
\end{cases}
\label{eq:condition-violation}
\end{equation}
The maximum of these condition-wise values gives~\eqref{eq:multi-capability}.

\subsection{Mechanism Indicators}
With tendon-constraint null-space basis $Z_{\mathcal A}$ and tip Jacobian $J_{\mathrm{tip}}$, the indicator
\begin{equation}
m(d)=\|J_{\mathrm{tip}}(d)Z_{\mathcal A}(d)\|_F^2
\label{eq:m0-indicator}
\end{equation}
measures local tip-motion freedom under the active tendon constraints. M0 evaluates its sensitivity to the selected structural parameters.

\subsection{Passive equilibrium sensitivity}
For active set $\mathcal A$, let $w_{\mathcal A}=(q,\mu_{\mathcal A})$ satisfy $F(w_{\mathcal A},d;f)=0$. When the active set is unchanged and $F_w$ is nonsingular, the implicit-function theorem gives $\partial w^\star/\partial d=-F_w^{-1}F_d$. A task metric follows by the chain rule. Loaded and unloaded equilibria are both differentiated for a response relative to an unloaded reference. Active-set transitions require directional or mode-specific analysis.

\subsection{Control-Optimized Sensitivity}
The epigraph form of~\eqref{eq:capability} is
\begin{equation}
\min_{U,\rho}\ \rho\quad\text{s.t.}\quad g_{\tau,i}(d,U)\leq\rho\ (\forall i),\quad U\in\mathcal U(d).
\label{eq:epigraph}
\end{equation}
Let $\lambda_i^\star$, $\alpha_j^\star$, and $\beta_j^\star$ denote the optimal multipliers for the task, lower-input, and upper-input constraints, respectively. The actuator bounds are $\ell_j(d)\leq U_j\leq h_j(d)$. Under the relevant local optimality, differentiability, constraint-qualification, and regularity conditions, the envelope relation gives
\begin{equation}
\begin{aligned}
\frac{\partial V_\tau}{\partial d_j}
={}&\sum_i\lambda_i^\star
\frac{\partial g_{\tau,i}}{\partial d_j}\\
&+\sum_\ell\alpha_\ell^\star
\frac{\partial \ell_\ell}{\partial d_j}
-\sum_\ell\beta_\ell^\star
\frac{\partial h_\ell}{\partial d_j}.
\end{aligned}
\label{eq:envelope-general}
\end{equation}
When actuator bounds are independent of the design, the last two terms vanish. The epigraph variable $\rho$ exposes the active task constraints and their multipliers. The relation applies to the regular local solution branch represented by these multipliers.

\subsection{Deterministic Capability Families}
\label{subsec:research-tools}
\begin{table*}[!t]
\centering
\caption{Capability-level interface between the Main Agent and deterministic robot tools.}
\label{tab:tool-action-taxonomy}
\papertablestyle
\begin{tabularx}{\textwidth}{@{}>{\itshape}p{0.19\textwidth}!{\color{TraceGray}\vrule width 0.25pt} X X@{}}
\toprule
\normalfont\textbf{Capability} & \textbf{Deterministic execution} & \textbf{Main Agent decision} \\
\midrule
Structure and modeling & Compile discrete morphology changes and continuous design candidates against the admissible design space & Select the physical mechanism, topology, and parameter group \\
Mathematical research & Execute registered M0--M3 analyses, validity checks, and numerical optimization & Select the research question, observables, analysis semantics, and budget \\
Control and trajectory & Solve task-specific actuation, initialization, trajectory, and feedback-control problems & Select the control formulation and decide when control should be re-optimized \\
Physical evaluation & Execute active dynamics or passive equilibrium and apply the frozen task evaluator & Select direct physical tests and interpret the resulting physical evidence \\
\bottomrule
\end{tabularx}
\end{table*}

Table~\ref{tab:tool-action-taxonomy} summarizes the capability interface. The research interface exposes a compact set of typed tool families. Each family has a deterministic execution contract, while the Main Agent chooses when and why it is used. The interface is organized around research actions rather than a flat list of software APIs.



The task, parameter, and observable registries provide deterministic bindings for the research interface. They resolve the frozen task contract, legal design variables and bounds, and registered physical quantities with their units and model scope. Before execution, the compiler checks identifiers, units, admissible ranges, and analysis applicability. Each capability returns typed evidence tied to its candidate, model, operating point, validity status, and computational cost. The Main Agent therefore selects research actions at the capability level without managing solver-specific APIs.
